\scanpage{85}
J'en arrive maintenant au cas \uline{B})
et je supposerai dorénavant, pour ce cas
\[
\boxed{2\cdot1_\Lambda\ne0}\tag*{(25)}
\]
de façon que \inserted{l'involution} $-\mathrm{id}_M$ opère librement sur
$\widetilde S_M=M^*$, sur $\widetilde A_M$ et sur $\widetilde F_M$ -- on
note \deleted{que} effet que
\[
\begin{aligned}
\{x,y,z\}\in\widetilde F_M\ \Longrightarrow\ &x,y,z,-x,-y,-z\ \text{sont}\\
&\text{tous distincts}.
\end{aligned}\tag*{(26)}
\]

On pose
\[
\left\{\begin{aligned}
S_M&=\widetilde S_M/\{\pm\mathrm{id}_M\}=M^*/\{\pm1\}\\
\deleted{A}F_M&=\deleted{\widetilde A}\widetilde F_M/\{\pm\mathrm{id}_M\}=
\end{aligned}\right.\tag*{(27)}
\]
et on a une bijection canonique
\[
F_M\xrightarrow{\sim}\left\{f\in\unsure{\mathfrak P|\mathbb P|\mathcal P}_3(S_M)\mid
\begin{aligned}&\exists\,\widetilde f\in\widetilde F_M,\ \text{avec}\\
&f=p(\widetilde f)\end{aligned}\right\}.\tag*{(28)}
\]
où
\[
\begin{aligned}p:\widetilde S_M&\longrightarrow S_M&\quad&\text{est l'application}\\
x&\longmapsto\dot x&&\text{canonique}.\end{aligned}\tag*{(29)}
\]

On vérifie \deleted{en effet} que si $f\in\mathfrak P_3(S_M)$
satisfait (28), alors il y a exactement deux
$\widetilde f$ tels que $f=p(\widetilde f)$, formant une orbite sous
$\{\pm\mathrm{id}_M\}$.
% archive sheet number: 84
